% A.81: full root and DE proof audits PASS; exact statement precision audited.
\subsection{Double covers with a non-CM elliptic Prym}

\begin{statement}{A.81}{Proved}{A general genus-two family and its wild torsion enlargements}
Let \(E:y^2=x^3-x\), and let
\(\psi:E'\to E\) be an actual isogeny of degree
\(n=2^a\ge2\), defined over a finite number field \(L\supset\mathbb Q(i)\).
Let \(q:Y\to E'\) be an actual connected double cover
of smooth proper curves, branched at exactly two distinct
points of \(\ker\psi\). Then \(Y\) has genus two.
Let \(\sigma\) be its deck involution and \(h\) its
hyperelliptic involution. The other elliptic quotient
\[
                  q_F:Y\longrightarrow F=Y/\langle\sigma h\rangle
                                                        \tag{A.81.1}
\]
is made an elliptic curve after a finite extension and a
choice of origin. Suppose that this actual elliptic curve
\(F\) is non-CM. Equivalently the elliptic Prym of \(q\)
is non-CM. Put
\[
       f=\psi q,\qquad S=q^{-1}(\ker\psi),\qquad X=Y\setminus S.
                                                        \tag{A.81.2}
\]
The actual degree-\(2n\) cover \(X\to E\setminus\{O\}\)
has \(2n-2\) cusps and two-group geometric Galois closure.
For every coefficient prime \(\ell\), original Part A packet
and permitted open \(\Delta_0\) preserving this cover and a
selected cusp lift \(P\),
\[
 \left(\pi_1(X_{\overline L})^{(\ell)}\right)^{s_b(\Delta_0)}
                              =I_{P,\ell}.             \tag{A.81.2a}
\]
Every nontrivial closed procyclic subgroup of
\(\pi_1(X_{\overline L})^{(\ell)}\) whose arithmetic
normalizer has open image in \(G_L\) is contained in a
conjugate of cusp inertia.
At odd coefficient primes the fixed assertion also holds
for every cusp packet over a finite number field of definition
and every open subgroup of its closed derived Galois group.

In particular, keep the actual curve \(Y\), map \(q_Y\)
and cyclic four-isogeny \(\eta\) of A.80. For every
\(r\ge0\), define
\[
 \psi_r=[2^r]\eta,\qquad
 X_r=Y\setminus q_Y^{-1}(\ker\psi_r).
                                                        \tag{A.81.3}
\]
Both assertions hold at every coefficient prime for these
actual covers. Their degrees and cusp counts are
\[
              \deg(X_r/E\setminus\{O\})=2\cdot4^{r+1},
              \qquad \#(Y\setminus X_r)=2\cdot4^{r+1}-2.
                                                        \tag{A.81.4}
\]
This is an unbounded family on one genus-two compactification.
It is not asserted cofinal among finite covers. The comparison
\(X_r\subset X_0\) is an open immersion with a changed
elliptic covering map, not a finite refinement or a fixed-lift
principle for arbitrary covers.
\end{statement}

\begin{proof}
\paragraph{Actual geometric quotients and closure.}
Make one finite extension defining the maps, branch points,
involutions and all cusps. Riemann--Hurwitz for the double
cover gives \(2g(Y)-2=2\), hence \(g(Y)=2\).
The hyperelliptic involution is central and acts as \(-1\)
on proper homology. The \(\sigma\)-invariant homology is
the dimension-two pullback from \(E'\); therefore the
\(\sigma h\)-invariant part is the complementary
dimension-two eigenspace. The actual quotient in (A.81.1)
has genus one. Its pullback and that of \(q\) give an
isogeny
\[
                       J_Y\sim E'\times F.             \tag{A.81.5}
\]
The image of \(q_F^*\) is the connected kernel of the
norm to \(E'\), up to isogeny. Thus the assumed Prym is
identified by an actual quotient map, not by an unspecified
isogeny factor. Translation of origins has no effect on
the proper homology or the non-CM hypothesis.

There is also an actual two-group closure. Over an algebraic
closure write \(k(Y)=k(E')(\sqrt a)\). For every
\(T\in\ker\psi\), adjoin the square root of the
translate \(a(P+T)\) in a common algebraic closure of
\(k(E')\). Their compositum is a connected field and
is Galois over \(k(E)\): the translations permute the
quadratic subextensions, and each quadratic extension is
already normal. Its group fits into an exact sequence
with elementary abelian two-group kernel and quotient
\(\ker\psi\), a two-group. Consequently the actual
normal closure of \(Y\to E\), a subextension of this
field, has two-group Galois group. No splitting of this
extension and no independence of the translated square
classes is needed.

The isogeny is unramified in characteristic zero. Both
branch points of \(q\) map to \(O\), so \(X\to E\setminus
\{O\}\) is finite etale. Above the \(n\) points of
\(\ker\psi\), the two branch points each have one
preimage and the remaining \(n-2\) points each have two.
This proves the exact cusp count \(2+2(n-2)=2n-2\).
In particular every selected cusp has at least one distinct
possible companion.

\paragraph{Every reduction hypothesis in the A.80 argument persists.}
At a place above two, \(E'\) has potential good
supersingular reduction, because it is isogenous to \(E\).
After finite extension all points of \(\ker\psi\)
specialize to the origin of its good special fibre.
A.78 applies to the actual one-point two-group closure
just constructed over \(E\). Norm maps make \(J_Y\)
a quotient of its Jacobian, so \(J_Y\) has potential
good reduction with special two-rank zero. Its proper
stable fibre therefore has exactly the alternatives used
in A.80: one smooth genus-two component, or two
supersingular elliptic components joined at a node.

The actual degree-two morphism \(q\) extends through
the relative Abel map. The sum of its positive component
degrees is two, by the constant degree of the pullback
of a degree-one line bundle on the good elliptic target.
These are precisely the actual degree configurations
\((2)\), \((1,1)\) and \((2,0)\) in A.80.
The involution has a nonzero anti-invariant proper
homology block, so cannot become the identity on every
positive-genus component. The two ramification cusps
are actual \(\sigma\)-fixed sections and both map to
the origin after reduction. None of these assertions
depends on the number of other cusps.

For completeness, the resulting path comparison is the
following. In the smooth genus-two case, the special
degree-two map is separable and its rank-zero
Deuring--Shafarevich formula gives a unique fixed point
over the origin. In degree \((1,1)\), the involution
exchanges the components and the ramification cusps
specialize to their node, which is the only point over
the origin on either component. All cusps thus lie in
a single rational marked insertion in these cases.
In degree \((2,0)\), the main elliptic component maps
radicially to \(E'_s\), while the involution on the
contracted supersingular elliptic tail has exactly one
fixed point, its node. If that node maps away from the
origin, all cusps lie over the single main point above
the origin. If it maps to the origin, both ramification
cusps specialize to the node. A selected cusp on the
tail may have no rational companion path, but a path
to either ramification cusp avoids the main elliptic
component with strictly positive distance.

In this last case the actual \(q_*\), together with
its polarization adjoint, identifies the global
\(V_\ell E'\) block with the main component and the
global \(V_\ell F\) block with the tail. The new
partial-distance proof A.80.10--A.80.16 therefore applies
with exactly these blocks and with all \(2n-3\)
independent cusp generators retained. It imposes no
bound on that number. It gives
\[
              \ker D\subset\mathbb L(T,V_\ell F,\Omega_{E'}).
                                                        \tag{A.81.6}
\]
Every packet term is allowed in that proof. The
coefficient of \(\Omega_{E'}\) in its own block
conjugator is the same strictly positive distance, so
the all-weight slot injection is unchanged.

The first global elliptic factor is CM, being isogenous
to \(E\), while \(F\) is non-CM by hypothesis.
The actual connected derived monodromy is therefore
\(\{1\}\times\mathrm{SL}(V_\ell F)\), as in
A.80.17. Its invariant tensors in (A.81.6) have even
weight and the corresponding scalar power of the
cyclotomic character. A.38, applied to the closure of
the actual fixed subgroup, proves fixed inertia at
every odd coefficient prime. The rational-path cases
are already covered by A.41. At coefficient two,
A.22 applies to the original Part A packets through
the actual two-group closure. A.59 supplies all the
normalizer assertions. Field restriction and geometric
paths are used exactly in their inner-invariant senses
in A.80; no rational lifting of an arbitrary Lie
element is asserted.

Finally, (A.81.3) preserves the two original branch
points and the actual non-CM quotient \(F\) of A.80.
The degree of \(\psi_r\) is \(4^{r+1}\), so the
general degree and cusp formulas give (A.81.4).
This is a fresh application to each actual enlarged
cover, rather than a transfer of fixedness through
the open immersion. No new numerical premise is
used: the exact non-CM and map certificates are those
already retained and independently rerun for A.80.
\end{proof}
